\documentclass[a4paper,11pt]{article}

\usepackage[brazil]{babel}
\usepackage[utf8]{inputenc}
\usepackage[T1]{fontenc}
\usepackage{lmodern}
\usepackage[margin=2.5cm]{geometry}
\usepackage{amsmath,amssymb}
\usepackage{booktabs}
\usepackage{float}
\usepackage{pgfplots}
\usepackage{tikz}
\usepackage{xcolor}
\usepackage{listings}
\usepackage{algorithm}
\usepackage{algpseudocode}
\usepackage[hidelinks]{hyperref}

\pgfplotsset{compat=1.18}
\usetikzlibrary{positioning}

\floatname{algorithm}{Algoritmo}
\renewcommand{\listalgorithmname}{Lista de algoritmos}
\algrenewcommand\algorithmicwhile{\textbf{enquanto}}
\algrenewcommand\algorithmicdo{\textbf{faça}}
\algrenewcommand\algorithmicfor{\textbf{para}}
\algrenewcommand\algorithmicforall{\textbf{para todo}}
\algrenewcommand\algorithmicif{\textbf{se}}
\algrenewcommand\algorithmicthen{\textbf{então}}
\algrenewcommand\algorithmicelse{\textbf{senão}}
\algrenewcommand\algorithmicreturn{\textbf{retorne}}
\algrenewcommand\algorithmicend{\textbf{fim}}

\lstset{basicstyle=\ttfamily\small, columns=fullflexible, frame=single, breaklines=true}

\title{\textbf{Trabalho Prático 01}\\[2pt]
\large Identificação de pontes (naïve e Tarjan) e caminhos eulerianos pelo método de Fleury}
\author{Bernardo Barbosa Heronville \and Luiz Felipe Volpi \and Pedro Henrique Lopes de Melo\\[6pt]
\small Teoria dos Grafos e Computabilidade -- Prof. Zenilton Kleber Gonçalves do Patrocínio Júnior\\
\small Ciência da Computação -- PUC Minas}
\date{Novembro de 2026}

\begin{document}

\maketitle

\begin{abstract}
Este trabalho implementa, em Java, dois métodos para identificar as pontes de um grafo simples
não-direcionado -- um método naïve, que remove cada aresta e testa a conectividade do grafo
resultante, e o método de Tarjan (1974) -- e usa cada um deles como subrotina do método de Fleury
para encontrar um caminho (ou ciclo) euleriano, ou determinar que ele não existe. Foram realizados
experimentos com grafos aleatórios eulerianos, semi-eulerianos e não eulerianos com 100, 1.000,
10.000 e 100.000 vértices. Os resultados confirmam o custo linear do método de Tarjan contra o custo
quadrático do método naïve para encontrar todas as pontes e mostram que o método de Fleury é
quadrático com qualquer das duas estratégias.
\end{abstract}

\section{Introdução}

Seja $G=(V,E)$ um grafo simples não-direcionado, com $n=|V|$ e $m=|E|$. Uma \emph{ponte} é uma
aresta cuja remoção aumenta o número de componentes conexos do grafo (no caso de um grafo conexo,
torna-o desconexo). Um \emph{caminho euleriano} é um caminho que passa por todas as arestas de $G$
exatamente uma vez; se ele começa e termina no mesmo vértice, é um \emph{ciclo euleriano}.

Pelo teorema de Euler, um grafo cujas arestas estão todas em um mesmo componente conexo
\begin{itemize}
  \item é \textbf{euleriano} (possui ciclo euleriano) se e somente se todos os vértices têm grau par;
  \item é \textbf{semi-euleriano} (possui caminho euleriano, mas não ciclo) se e somente se possui
        exatamente dois vértices de grau ímpar -- e todo caminho euleriano começa em um deles e termina no outro;
\end{itemize}
e, em qualquer outro caso, o grafo é \textbf{não euleriano}.

O método de Fleury~\cite{fleury} constrói o caminho atravessando, a cada passo, uma aresta incidente ao vértice
atual e apagando-a, com a única restrição de \emph{não atravessar uma ponte do grafo restante, a menos
que não haja alternativa}. Por isso ele depende de um procedimento de identificação de pontes, que é o
foco deste trabalho.

O código-fonte está na pasta \texttt{codigo/} e este relatório (com seus fontes em \TeX) na pasta
\texttt{relatorio/}. O Apêndice~\ref{ap:uso} descreve como compilar e executar.

\section{Implementação}

\subsection{Representação do grafo}

A classe \texttt{Grafo} armazena os vértices numerados de $1$ a $n$ e as arestas numeradas de $0$
a $m-1$, cada uma com suas extremidades \texttt{origem[e]} e \texttt{destino[e]}. A lista de
adjacência de cada vértice guarda \emph{identificadores de arestas} e cada aresta guarda a sua posição
na lista de cada extremidade. Com isso:

\begin{itemize}
  \item \textbf{remover} uma aresta custa $O(1)$: ela é trocada com a última aresta de cada lista e o
        grau dos dois vértices é decrementado;
  \item \textbf{reinserir} uma aresta (necessário no teste naïve, que remove e devolve) custa $O(1)$;
  \item percorrer a vizinhança de $v$ visita apenas arestas ainda presentes, em $O(\deg(v))$.
\end{itemize}

Essa representação evita o custo de pular arestas já apagadas, o que seria importante no Fleury,
em que todas as arestas acabam sendo removidas. A leitura de arquivo descarta laços e arestas
repetidas, já que o enunciado trata de grafos simples. Todas as buscas são \emph{iterativas}, com pilha
ou fila explícita, para que grafos com 100.000 vértices não estourem a pilha de chamadas da JVM.
Os vetores auxiliares das buscas são alocados uma única vez por grafo; em vez de limpá-los a cada
busca, usa-se uma marcação por ``rodada'' (um vértice está visitado se \texttt{visitado[v]} é igual ao
número da rodada atual).

\subsection{Método naïve}

O método naïve segue a definição de ponte: para cada aresta $e$, remove-se $e$, conta-se o número de
componentes conexos com uma busca em largura e devolve-se $e$ ao grafo. A aresta é ponte se o número de
componentes aumentou (Algoritmo~\ref{alg:naive}). Cada teste custa $O(n+m)$, logo encontrar todas as
pontes custa $O(m(n+m))$.

\begin{algorithm}[H]
\caption{Pontes pelo método naïve}\label{alg:naive}
\begin{algorithmic}[1]
\State $c_0 \gets$ número de componentes de $G$ \Comment{BFS, $O(n+m)$}
\State $P \gets \emptyset$
\ForAll{aresta $e \in E$}
  \State remova $e$ de $G$
  \If{número de componentes de $G > c_0$} $P \gets P \cup \{e\}$ \EndIf
  \State reinsira $e$ em $G$
\EndFor
\State \Return $P$
\end{algorithmic}
\end{algorithm}

Dentro do Fleury, não é preciso conhecer todas as pontes, apenas saber se uma aresta candidata
$(u,v)$ é ponte. O teste naïve de uma única aresta (\texttt{ehPonteNaive}) remove $(u,v)$, faz uma busca
em largura a partir de $u$ e verifica se $v$ continua alcançável -- isto é, se o componente de $u$
continua conexo. Se $v$ não é alcançado, $(u,v)$ é ponte. O custo é $O(n+m)$ por teste.

\subsection{Método de Tarjan (1974)}

Implementamos o algoritmo descrito por Tarjan em \emph{A note on finding the bridges of a graph}
\cite{tarjan1974}. A ideia é que uma aresta de árvore $(\mathrm{pai}(w), w)$ de uma árvore geradora $T$
só pode ser ponte se nenhuma aresta fora de $T$ ligar a subárvore de $w$ a um vértice fora dela (arestas
fora da árvore geradora nunca são pontes, pois fecham um ciclo com $T$). Para cada componente:

\begin{enumerate}
  \item Construir uma árvore geradora $T$ e numerar os vértices em pré-ordem de $T$, $\mathrm{pre}(v)$.
        Usamos uma busca em profundidade iterativa; com essa numeração os descendentes de $w$ são
        exatamente os vértices com números em $[\mathrm{pre}(w),\ \mathrm{pre}(w)+ND(w)-1]$.
  \item Calcular $ND(v)$, o número de descendentes de $v$ em $T$ (incluindo o próprio $v$).
  \item Calcular, em pós-ordem,
  \begin{align*}
    L(v) &= \min\big(\{\mathrm{pre}(v)\} \cup \{L(w) : w \text{ filho de } v\} \cup \{\mathrm{pre}(w) : (v,w)\in E\setminus T\}\big),\\
    H(v) &= \max\big(\{\mathrm{pre}(v)\} \cup \{H(w) : w \text{ filho de } v\} \cup \{\mathrm{pre}(w) : (v,w)\in E\setminus T\}\big),
  \end{align*}
  ou seja, o menor e o maior número de pré-ordem alcançáveis a partir da subárvore de $v$ usando no
  máximo uma aresta fora de $T$.
  \item A aresta de árvore $(\mathrm{pai}(w), w)$ é ponte se e somente se
  \[ L(w) = \mathrm{pre}(w) \quad\text{e}\quad H(w) < \mathrm{pre}(w) + ND(w). \]
\end{enumerate}

Como a pré-ordem coloca todo vértice antes dos seus descendentes, percorrer o vetor da pré-ordem
\emph{de trás para frente} processa cada vértice depois de todos os seus filhos, o que dá a pós-ordem
dos passos 2 e 3 sem recursão: ao terminar um vértice $v$, testa-se a condição do passo 4 para a aresta
que o liga ao pai e propagam-se $ND(v)$, $L(v)$ e $H(v)$ para o pai. Cada vértice e cada aresta são
examinados um número constante de vezes, então o custo é $O(n+m)$. O Algoritmo~\ref{alg:tarjan}
resume a implementação (\texttt{Pontes.tarjanComponente}).

\begin{algorithm}[H]
\caption{Pontes pelo método de Tarjan (1974), para o componente de $r$}\label{alg:tarjan}
\begin{algorithmic}[1]
\State faça uma DFS iterativa a partir de $r$, registrando a pré-ordem $\mathit{ordem}[0..k-1]$,
       $\mathrm{pre}(v)$ e a aresta-pai de cada vértice
\State $ND(v)\gets 1$; $L(v)\gets H(v)\gets\mathrm{pre}(v)$ para todo $v$ visitado
\For{$i \gets k-1$ \textbf{até} $0$}
  \State $v \gets \mathit{ordem}[i]$
  \ForAll{aresta $(v,w)$ que não pertence a $T$}
     \State $L(v)\gets\min(L(v),\mathrm{pre}(w))$; \quad $H(v)\gets\max(H(v),\mathrm{pre}(w))$
  \EndFor
  \If{$v \neq r$}
     \State $p \gets \mathrm{pai}(v)$
     \If{$L(v)=\mathrm{pre}(v)$ \textbf{e} $H(v)<\mathrm{pre}(v)+ND(v)$} marque $(p,v)$ como ponte \EndIf
     \State $ND(p)\gets ND(p)+ND(v)$; \quad $L(p)\gets\min(L(p),L(v))$; \quad $H(p)\gets\max(H(p),H(v))$
  \EndIf
\EndFor
\end{algorithmic}
\end{algorithm}

Para encontrar todas as pontes de um grafo possivelmente desconexo, o procedimento é repetido a partir
de cada vértice ainda não visitado (uma árvore por componente, formando a floresta geradora do artigo).

\subsection{Classificação e método de Fleury}

A função \texttt{Fleury.classificar} conta os vértices de grau ímpar e verifica, por uma busca em
largura, se todos os vértices com grau positivo estão no mesmo componente (vértices isolados são
ignorados, pois o caminho só precisa cobrir as arestas). Isso custa $O(n+m)$ e decide a existência
do caminho: se o grafo é não euleriano, o programa informa que não há caminho euleriano e Fleury não é
executado.

Caso exista caminho, ele começa em um vértice de grau ímpar (semi-euleriano) ou em qualquer vértice
com arestas (euleriano), e o Algoritmo~\ref{alg:fleury} é executado sobre uma cópia do grafo.

\begin{algorithm}[H]
\caption{Método de Fleury}\label{alg:fleury}
\begin{algorithmic}[1]
\State $u \gets$ vértice inicial; $C \gets (u)$
\While{$u$ tem arestas}
  \If{$\deg(u) = 1$}
    \State $e \gets$ a única aresta de $u$ \Comment{não há alternativa}
  \Else
    \State $e \gets$ uma aresta de $u$ que \textbf{não} é ponte do grafo atual
  \EndIf
  \State $w \gets$ outra extremidade de $e$; remova $e$; acrescente $w$ a $C$; $u \gets w$
\EndWhile
\State \Return $C$
\end{algorithmic}
\end{algorithm}

A linha 6 é onde entram as duas estratégias:

\begin{description}
  \item[Fleury + naïve:] as arestas de $u$ são testadas uma a uma com \texttt{ehPonteNaive} até
        encontrar a primeira que não é ponte.
  \item[Fleury + Tarjan:] executa-se o método de Tarjan no componente de $u$ (que, durante o Fleury,
        contém todas as arestas restantes) e escolhe-se a primeira aresta de $u$ não marcada como ponte.
\end{description}

Quando $\deg(u)\ge 2$ em um grafo que ainda possui caminho euleriano a partir de $u$, no máximo uma
aresta de $u$ é ponte; logo a estratégia naïve faz no máximo dois testes por passo e quase sempre faz
apenas um. Em ambos os casos o custo de um passo é $O(n+m)$, e o Fleury completo custa $O(m(n+m))$.
Ao final, a função \texttt{Fleury.valido} confere que o caminho tem $m+1$ vértices e usa cada aresta
exatamente uma vez; essa verificação é feita em todos os experimentos.

\subsection{Geração dos grafos aleatórios}

A classe \texttt{GeradorGrafos} produz grafos simples e conexos com grau médio aproximadamente igual a
4 ($m\approx 2n$):

\begin{description}
  \item[Euleriano:] começa com um ciclo hamiltoniano sobre uma permutação aleatória dos vértices
        (o que garante conexidade e graus pares) e acrescenta ciclos aleatórios de 3 a 8 vértices cujas
        arestas ainda não existem. Cada ciclo soma 2 ao grau de seus vértices, mantendo todos os graus pares.
  \item[Semi-euleriano:] um grafo euleriano do qual se remove uma aresta aleatória. O grafo continua
        conexo (o ciclo hamiltoniano sem uma aresta ainda é um caminho hamiltoniano) e passa a ter
        exatamente dois vértices de grau ímpar.
  \item[Não euleriano:] um grafo euleriano ao qual se acrescentam $k=\max(2, n/50)$ arestas novas com
        extremidades todas distintas, criando $2k\ge 4$ vértices de grau ímpar.
\end{description}

Ao final, a ordem das arestas é embaralhada para que a ordem de leitura não favoreça nenhum método.
Todo grafo gerado é classificado novamente pelo programa antes das medições, para garantir que a classe
está correta.

\subsection{Testes de corretude}

A classe \texttt{Testes} verifica: (i) o grafo da Figura~1 do enunciado, em que ambos os métodos
encontram apenas a ponte $F$--$G$ e o grafo é classificado como não euleriano (os vértices $A$, $E$, $F$
e $G$ têm grau ímpar); (ii) 3.000 grafos aleatórios pequenos, construídos a partir de florestas e
portanto cheios de pontes, nos quais os métodos naïve e de Tarjan devem devolver exatamente o mesmo
conjunto de pontes; (iii) 900 grafos gerados das três classes, nos quais o Fleury com as duas estratégias
deve devolver um caminho válido (fechado, no caso euleriano) ou nenhum caminho (não euleriano).
Todos os testes passam.

\section{Metodologia experimental}

Para cada $n\in\{100;\ 1.000;\ 10.000;\ 100.000\}$ e cada classe de grafo, foram gerados grafos
aleatórios com semente fixa (o que torna os experimentos reproduzíveis) e medidos, com
\texttt{System.nanoTime()}, os tempos de quatro tarefas:

\begin{itemize}
  \item \textbf{Pontes naïve} e \textbf{Pontes Tarjan}: encontrar todas as pontes do grafo;
  \item \textbf{Fleury + naïve} e \textbf{Fleury + Tarjan}: classificar o grafo e, se possível,
        encontrar o caminho euleriano pelo método de Fleury.
\end{itemize}

Foram usados 10 grafos por classe para $n=100$ e $n=1.000$, 3 para $n=10.000$ e 1 para $n=100.000$
(cada execução de um método quadrático nesse tamanho leva vários minutos). Antes das medições, todas as
tarefas são executadas várias vezes em grafos pequenos para que o compilador JIT da JVM já tenha
compilado o código. Cada medição tem limite de 30 minutos; um método que estoure o limite não é
executado nos tamanhos maiores. Em todas as execuções os conjuntos de pontes dos dois métodos foram
comparados e os caminhos encontrados foram validados.

\textbf{Ambiente:} Apple M2 com 8~GB de memória, macOS, OpenJDK 24.0.2, execução com um único thread.

\section{Resultados}

A Tabela~\ref{tab:resultados} mostra os tempos médios e as
Figuras~\ref{fig:euleriano}, \ref{fig:semi-euleriano} e \ref{fig:nao-euleriano} apresentam os mesmos
dados em escala log-log. Nessa escala, uma reta de inclinação 1 indica crescimento linear e uma reta
de inclinação 2, crescimento quadrático.

\begin{table}[H]
\centering
\small
\setlength{\tabcolsep}{4pt}
\caption{Tempo médio (ms) de cada estratégia. $\bar m$ é o número médio de arestas e $r$ o número de grafos.}
\label{tab:resultados}
\begin{tabular}{llrrrrrr}
\toprule
Classe & $n$ & $\bar m$ & $r$ & Pontes naïve & Pontes Tarjan & Fleury+naïve & Fleury+Tarjan \\
\midrule
Euleriano & 100 & 202 & 10 & 0,268 & 0,018 & 0,383 & 0,470 \\
 & 1.000 & 2.002 & 10 & 15,80 & 0,124 & 20,55 & 39,56 \\
 & 10.000 & 20.004 & 3 & 6.757 & 1,29 & 2.568 & 7.391 \\
 & 100.000 & 200.000 & 1 & 832.585 & 19,87 & 294.380 & 889.606 \\
\midrule
Semi-euleriano & 100 & 201 & 10 & 0,254 & 0,025 & 1,77 & 0,395 \\
 & 1.000 & 2.001 & 10 & 16,05 & 0,119 & 20,12 & 39,13 \\
 & 10.000 & 20.002 & 3 & 6.816 & 1,24 & 2.562 & 7.348 \\
 & 100.000 & 200.001 & 1 & 822.728 & 17,18 & 285.468 & 882.548 \\
\midrule
Não euleriano & 100 & 204 & 10 & 0,266 & 0,025 & 0,0005 & 0,0003 \\
 & 1.000 & 2.022 & 10 & 16,41 & 0,104 & 0,0012 & 0,0009 \\
 & 10.000 & 20.201 & 3 & 6.884 & 1,25 & 0,0086 & 0,0065 \\
 & 100.000 & 202.001 & 1 & 837.576 & 16,89 & 0,103 & 0,083 \\
\bottomrule
\end{tabular}
\end{table}

\begin{figure}[H]
\centering
\begin{tikzpicture}
\begin{loglogaxis}[width=0.85\textwidth, height=6.2cm, xlabel={Número de vértices $n$}, ylabel={Tempo médio (ms)}, grid=major, legend pos=north west, legend style={font=\footnotesize}, xtick={100,1000,10000,100000}]
\addplot[color=red, mark=square*] coordinates {(100,0.26809) (1000,15.8007) (10000,6756.98) (100000,832585)};
\addlegendentry{Pontes -- naïve}
\addplot[color=blue, mark=*] coordinates {(100,0.01849) (1000,0.12373) (10000,1.28687) (100000,19.8658)};
\addlegendentry{Pontes -- Tarjan}
\addplot[color=orange, mark=triangle*, dashed] coordinates {(100,0.38273) (1000,20.5541) (10000,2568.21) (100000,294380)};
\addlegendentry{Fleury + naïve}
\addplot[color=teal, mark=diamond*, dashed] coordinates {(100,0.47039) (1000,39.556) (10000,7390.7) (100000,889606)};
\addlegendentry{Fleury + Tarjan}
\end{loglogaxis}
\end{tikzpicture}
\caption{Tempo médio em escala log-log -- grafos da classe euleriano.}
\label{fig:euleriano}
\end{figure}

\begin{figure}[H]
\centering
\begin{tikzpicture}
\begin{loglogaxis}[width=0.85\textwidth, height=6.2cm, xlabel={Número de vértices $n$}, ylabel={Tempo médio (ms)}, grid=major, legend pos=north west, legend style={font=\footnotesize}, xtick={100,1000,10000,100000}]
\addplot[color=red, mark=square*] coordinates {(100,0.25368) (1000,16.0522) (10000,6816.04) (100000,822728)};
\addlegendentry{Pontes -- naïve}
\addplot[color=blue, mark=*] coordinates {(100,0.02549) (1000,0.11882) (10000,1.244) (100000,17.183)};
\addlegendentry{Pontes -- Tarjan}
\addplot[color=orange, mark=triangle*, dashed] coordinates {(100,1.76888) (1000,20.125) (10000,2561.79) (100000,285468)};
\addlegendentry{Fleury + naïve}
\addplot[color=teal, mark=diamond*, dashed] coordinates {(100,0.39473) (1000,39.1314) (10000,7347.59) (100000,882548)};
\addlegendentry{Fleury + Tarjan}
\end{loglogaxis}
\end{tikzpicture}
\caption{Tempo médio em escala log-log -- grafos da classe semi-euleriano.}
\label{fig:semi-euleriano}
\end{figure}

\begin{figure}[H]
\centering
\begin{tikzpicture}
\begin{loglogaxis}[width=0.85\textwidth, height=6.2cm, xlabel={Número de vértices $n$}, ylabel={Tempo médio (ms)}, grid=major, legend pos=north west, legend style={font=\footnotesize}, xtick={100,1000,10000,100000}]
\addplot[color=red, mark=square*] coordinates {(100,0.26641) (1000,16.407) (10000,6883.65) (100000,837576)};
\addlegendentry{Pontes -- naïve}
\addplot[color=blue, mark=*] coordinates {(100,0.02527) (1000,0.10424) (10000,1.25067) (100000,16.8857)};
\addlegendentry{Pontes -- Tarjan}
\addplot[color=orange, mark=triangle*, dashed] coordinates {(100,0.00046) (1000,0.00121) (10000,0.00863333) (100000,0.1032)};
\addlegendentry{Fleury + naïve}
\addplot[color=teal, mark=diamond*, dashed] coordinates {(100,0.00027) (1000,0.00091) (10000,0.0065) (100000,0.0826)};
\addlegendentry{Fleury + Tarjan}
\end{loglogaxis}
\end{tikzpicture}
\caption{Tempo médio em escala log-log -- grafos da classe não euleriano.}
\label{fig:nao-euleriano}
\end{figure}


\section{Discussão}

Todas as medições terminaram dentro do limite de 30 minutos. Em todas as execuções os dois métodos
encontraram o mesmo conjunto de pontes e todos os caminhos eulerianos encontrados foram validados.
Os grafos eulerianos e não eulerianos gerados contêm um ciclo hamiltoniano e, portanto, não têm pontes.
Nos semi-eulerianos, a aresta removida pode deixar um vértice de grau 1; se ele estava no fim de uma
cadeia de vértices de grau 2 do ciclo hamiltoniano, todas as arestas dessa cadeia viram pontes. Dos 24
grafos semi-eulerianos usados, 6 tinham de 1 a 3 pontes. O número de pontes não altera o
custo dos métodos, que examinam todas as arestas de qualquer forma. A presença de pontes também foi
exercitada nos testes de corretude e, durante o Fleury, nos grafos intermediários, que passam a ter
pontes à medida que as arestas são apagadas.

\subsection{Identificação de todas as pontes}

A diferença entre os dois métodos é a esperada pela análise assintótica. Com $n=100.000$ ($m\approx
200.000$), o método de Tarjan encontrou todas as pontes em cerca de \textbf{17 a 20~ms}, enquanto o
método naïve levou cerca de \textbf{830 segundos} (quase 14 minutos) -- aproximadamente
\textbf{45.000 vezes} mais. A inclinação das curvas no gráfico log-log mede o expoente de crescimento
(Tabela~\ref{tab:inclinacao}):

\begin{table}[H]
\centering
\small
\caption{Inclinação $\log_{10}(t_{10n}/t_n)$ entre tamanhos consecutivos (grafos eulerianos). Valores
próximos de 1 indicam crescimento linear e próximos de 2, quadrático.}
\label{tab:inclinacao}
\begin{tabular}{lccc}
\toprule
Medida & $100\to1.000$ & $1.000\to10.000$ & $10.000\to100.000$ \\
\midrule
Pontes naïve   & 1,77 & 2,63 & 2,09 \\
Pontes Tarjan  & 0,83 & 1,02 & 1,19 \\
Fleury + naïve & 1,73 & 2,10 & 2,06 \\
Fleury + Tarjan& 1,92 & 2,27 & 2,08 \\
\bottomrule
\end{tabular}
\end{table}

\begin{itemize}
  \item \textbf{Tarjan} cresce linearmente ($O(n+m)$): a partir de 1.000 vértices, multiplicar o tamanho
        por 10 multiplica o tempo por 10 a 15. O pequeno excesso acima de 1 entre 10.000 e 100.000 vértices provavelmente se deve
        à memória cache: os vetores de um grafo com 100.000 vértices não cabem mais nas caches do
        processador e os acessos da busca, que saltam entre vértices aleatórios, ficam mais caros.
  \item \textbf{Naïve} cresce quadraticamente ($O(m(n+m))$): são $m$ buscas em largura, cada uma de
        custo $O(n+m)$. A inclinação 2,63 entre 1.000 e 10.000 vértices provavelmente também é
        efeito de cache -- um grafo com 1.000 vértices cabe inteiro na cache e cada busca é muito rápida;
        com 10.000 já não cabe (não medimos isso diretamente). Entre 10.000 e 100.000 a inclinação volta a ser $\approx 2$.
  \item Os tempos praticamente não dependem da classe do grafo (euleriano, semi ou não euleriano),
        pois os três tipos de grafo têm o mesmo tamanho e a mesma estrutura geral.
\end{itemize}

\subsection{Método de Fleury}

Para grafos eulerianos e semi-eulerianos, o Fleury é \textbf{quadrático com as duas estratégias}
(inclinação $\approx 2$): ele atravessa $m$ arestas e, a cada passo em que o vértice atual tem grau
maior que 1, faz uma identificação de ponte de custo $O(n+m)$. Com 100.000 vértices, o Fleury levou
cerca de \textbf{4,8~minutos} com o teste naïve e \textbf{14,8~minutos} com Tarjan.

Pode parecer contraditório que, dentro do Fleury, a estratégia naïve seja cerca de \textbf{três vezes
mais rápida} que Tarjan, mas a explicação é que as duas fazem trabalhos diferentes a cada passo:

\begin{itemize}
  \item a estratégia \emph{naïve} só precisa saber se \emph{uma} aresta candidata é ponte. Como no
        máximo uma aresta de $u$ pode ser ponte, em geral a primeira candidata já não é ponte e o passo
        custa uma única busca em largura -- um percurso simples sobre o componente;
  \item a estratégia \emph{Tarjan} calcula \emph{todas} as pontes do componente a cada passo: uma busca em
        profundidade com pilha explícita e, depois, uma segunda passada sobre todas as arestas para
        calcular $ND$, $L$ e $H$. O custo também é $O(n+m)$, mas com constante maior, e quase toda essa
        informação é descartada no passo seguinte, pois o grafo muda.
\end{itemize}

Ou seja, a vantagem assintótica de Tarjan aparece quando se quer \emph{todas} as pontes de um grafo fixo
(o que ele faz em uma passada, contra $m$ passadas do naïve). No Fleury, em que a pergunta é sobre uma
única aresta e o grafo muda a cada passo, as duas estratégias ficam na mesma ordem $O(m(n+m))$ e a de
menor constante vence.

Para grafos \textbf{não eulerianos}, o Fleury nem chega a ser executado: a classificação (contagem de
vértices de grau ímpar) detecta que o caminho não existe. Esse tempo é linear e muito pequeno
(0,1~ms com 100.000 vértices), igual para as duas estratégias.

Uma observação sobre a média: em grafos semi-eulerianos com $n=100$, uma das dez execuções do Fleury
com o teste naïve levou 15,6~ms, enquanto as demais ficaram em torno de 0,23~ms, o que eleva a média
para 1,77~ms. Essa execução isolada coincidiu provavelmente com uma pausa do coletor de lixo ou uma
recompilação da JVM, e não com o algoritmo. Em tempos tão pequenos esse ruído é comum; a partir de
1.000 vértices ele deixa de ser relevante.


\section{Conclusão}

Foram implementados os dois métodos de identificação de pontes pedidos -- o naïve, que testa a
conectividade após a remoção de cada aresta, e o de Tarjan (1974), baseado em uma árvore geradora
numerada em pré-ordem e nos valores $ND$, $L$ e $H$ -- e o método de Fleury usando cada um deles.
O programa classifica qualquer grafo como euleriano, semi-euleriano ou não euleriano e, quando existe,
devolve um caminho euleriano validado.

Os experimentos confirmam a análise assintótica: para encontrar todas as pontes, Tarjan é linear e o
método naïve é quadrático, com diferença de várias ordens de grandeza a partir de 10.000 vértices.
Já no método de Fleury as duas estratégias são quadráticas, porque a identificação de pontes é refeita
a cada uma das $m$ arestas atravessadas; nesse uso o teste naïve de uma única aresta mostrou-se até
mais barato que recalcular todas as pontes com Tarjan. Melhorias possíveis seriam usar estruturas de
conectividade dinâmica para o Fleury ou simplesmente o algoritmo de Hierholzer, que encontra o caminho
euleriano em $O(n+m)$.

\begin{thebibliography}{9}
\bibitem{tarjan1974}
TARJAN, R. E. A note on finding the bridges of a graph. \emph{Information Processing Letters},
v.~2, n.~6, p.~160--161, 1974. \url{https://doi.org/10.1016/0020-0190(74)90003-9}.

\bibitem{fleury}
FLEURY, M. Deux problèmes de géométrie de situation. \emph{Journal de mathématiques élémentaires},
v.~2, p.~257--261, 1883.
\end{thebibliography}

\appendix
\section{Como compilar e executar}\label{ap:uso}

\begin{lstlisting}
cd codigo
javac *.java

java Testes                              # testes de corretude
java Main ../exemplos/figura1.txt        # pontes + caminho euleriano de um arquivo
java GeradorGrafos 1000 0 42 g.txt       # gera grafo (0=euleriano, 1=semi, 2=nao)
java Experimentos 1800 4 ../resultados/tempos.csv   # experimentos (limite 1800 s, grau medio 4)

cd ../relatorio
python3 gerar_resultados.py              # tabela e graficos a partir do CSV
tectonic relatorio.tex                   # ou pdflatex relatorio.tex
\end{lstlisting}

O arquivo de entrada tem o mesmo formato usado nas tarefas da disciplina: a primeira linha contém
$n$ e $m$ e cada uma das $m$ linhas seguintes contém uma aresta $u\ v$, com vértices de $1$ a $n$.
Para grafos com mais de 200 arestas, o \texttt{Main} grava o caminho em \texttt{caminho-naive.txt} e
\texttt{caminho-tarjan.txt} em vez de imprimi-lo.

\end{document}
